\chapter{Engenharia de Software e Modernização}

\section{Engenharia de Requisitos com Sommerville}

Sommerville~\cite{sommerville2011} estrutura a engenharia de requisitos em quatro fases
interdependentes: elicitação, análise, especificação e validação. A reestruturação do
sistema em três módulos — \texttt{models.py}, \texttt{components.py} e \texttt{app.py}
— evidencia como cada fase foi aplicada para preservar as regras de negócio do legado
Vaadin/Java sem reaproveitar seu código-fonte.

\begin{description}

\item[a)] \textbf{Elicitação}: As entradas críticas do domínio de regressão espectral foram
identificadas diretamente no sistema Java original: as matrizes \texttt{Xcal} e
\texttt{Xteste}, os vetores \texttt{Ycal} e \texttt{Yteste}, e os hiperparâmetros
\textit{m} (número de bags) e \textit{k} (fração de subagging). O \texttt{app.py}
captura esses insumos por meio de \texttt{st.file\_uploader} e \texttt{st.text\_input},
incluindo suporte a decimal com vírgula — particularidade do padrão brasileiro — tratada
nas funções \texttt{parse\_matrix()} e \texttt{parse\_vector()} do \texttt{models.py}.

\item[b)] \textbf{Análise}: A análise revelou que a lógica do \texttt{RegressionService.java}
podia ser decomposta em responsabilidades distintas: parse de arquivos, algoritmos de
regressão, métricas de avaliação e otimização por validação cruzada. Essa decomposição
originou a estrutura modular adotada, onde \texttt{models.py} concentra toda a
inteligência matemática, sem qualquer dependência de interface.

\item[c)] \textbf{Especificação}: Os contratos de entrada e saída foram formalizados com
\textit{type hints} explícitos em todas as funções do \texttt{models.py}:
\texttt{parse\_matrix(content: bytes) -> np.ndarray},
\texttt{perform\_rlm(Xcal, Ycal, Xteste: np.ndarray) -> np.ndarray},
\texttt{calc\_metrics(actual, predicted) -> tuple[float, float, float]}.
A função \texttt{run\_full\_analysis()} retorna um dicionário tipado com chaves
\texttt{log}, \texttt{optimal}, \texttt{metrics} e \texttt{preds}, constituindo o
contrato formal entre a camada de negócio e a camada de apresentação.

\item[d)] \textbf{Validação}: A validação cruzada 5-fold implementada em
\texttt{run\_cross\_validation()} opera como oráculo de equivalência, permitindo comparar
o desempenho do sistema Python com os resultados históricos do sistema Java. O log
estruturado gerado por \texttt{run\_full\_analysis()} documenta cada etapa — dimensões
detectadas, RMSE-CV por combinação de hiperparâmetros, métricas finais e tempo de
execução — possibilitando auditoria científica e rastreabilidade completa.

\end{description}

% ─────────────────────────────────────────────────────────────────────────────
\section{Especificação de Requisitos Funcionais e Não Funcionais}
\label{sec:requisitos}

A fase de Especificação de Sommerville~\cite{sommerville2011}, descrita no item (c) da
seção anterior, é formalizada nesta seção por meio da enumeração dos requisitos
funcionais (RF) e não funcionais (RNF) efetivamente implementados no sistema
Python/Streamlit. Os requisitos foram extraídos diretamente do código-fonte de
\texttt{models.py}, \texttt{components.py} e \texttt{app.py}, e não de funcionalidades
planejadas ou não concluídas, de modo a refletir com precisão o estado real do sistema
entregue neste TCC~II.

\subsection{Requisitos Funcionais}

O Quadro~\ref{qua:rf} apresenta os requisitos funcionais do sistema, isto é, as
funcionalidades que o usuário pode efetivamente executar por meio da interface
Streamlit.

\begin{quadro}[H]
\centering
\caption{Requisitos funcionais do sistema}
\label{qua:rf}
\begin{tabularx}{\textwidth}{|l|X|}
\hline
\textbf{ID} & \textbf{Descrição} \\ \hline
RF01 & Permitir o \textit{upload} das matrizes \texttt{Xcal}, \texttt{Xteste} e dos
vetores \texttt{Ycal}, \texttt{Yteste} nos formatos \texttt{.txt}, \texttt{.csv} e
\texttt{.xlsx}. \\ \hline
RF02 & Interpretar arquivos com separador decimal brasileiro (vírgula) e separador de
coluna \texttt{;} ou \texttt{,}, normalizando-os para o formato numérico do NumPy. \\
\hline
RF03 & Executar a Regressão Linear Múltipla (RLM) simples via \texttt{np.linalg.lstsq}.
\\ \hline
RF04 & Executar RLM com \textit{Bagging} (reamostragem \textit{bootstrap} com reposição)
para um conjunto de valores de $m$ informados pelo usuário. \\ \hline
RF05 & Executar RLM com \textit{Subagging} (reamostragem sem reposição) para conjuntos de
valores de $m$ e $k$ informados pelo usuário. \\ \hline
RF06 & Otimizar os hiperparâmetros $m$ (\textit{Bagging}) e $k$, $m$ (\textit{Subagging})
por \textit{grid search} com validação cruzada de 5 \textit{folds}. \\ \hline
RF07 & Calcular as métricas RMSE, MAE e R$^2$ para os três modelos (RLM, \textit{Bagging}
e \textit{Subagging}) sobre o conjunto de teste. \\ \hline
RF08 & Exibir um log estruturado de execução, contendo dimensões detectadas, RMSE-CV por
combinação de hiperparâmetros e tempo de execução por modelo. \\ \hline
RF09 & Gerar gráficos interativos (Plotly) de valores reais \textit{versus} preditos por
modelo e de dispersão real \textit{versus} predito com os três modelos sobrepostos. \\
\hline
RF10 & Exportar a tabela comparativa de métricas em \texttt{.csv} e o log de execução em
\texttt{.txt}. \\ \hline
RF11 & Exibir a interface e os rótulos dos gráficos em português ou inglês, conforme
seleção do usuário. \\ \hline
RF12 & Exibir barra de progresso durante a execução do \textit{grid search}, atualizada
por \textit{callback} a cada combinação de hiperparâmetros avaliada. \\ \hline
\end{tabularx}
\begin{tablenotes}[flushleft]
\footnotesize
\item Fonte:~\oautor
\end{tablenotes}
\end{quadro}

\subsection{Requisitos Não Funcionais}

O Quadro~\ref{qua:rnf} apresenta os requisitos não funcionais, relativos a atributos de
qualidade do sistema — desempenho, manutenibilidade, testabilidade e confiabilidade —
verificáveis independentemente das funcionalidades de negócio listadas no
Quadro~\ref{qua:rf}.

\begin{quadro}[H]
\centering
\caption{Requisitos não funcionais do sistema}
\label{qua:rnf}
\begin{tabularx}{\textwidth}{|l|X|}
\hline
\textbf{ID} & \textbf{Descrição} \\ \hline
RNF01 & \textbf{Modularidade}: o sistema deve ser estruturado em módulos com acoplamento
zero — \texttt{models.py} não importa \texttt{streamlit} nem \texttt{plotly};
\texttt{components.py} não importa \texttt{streamlit} nem \texttt{numpy}. \\ \hline
RNF02 & \textbf{Testabilidade}: a lógica de negócio deve ser testável de forma isolada,
sem instanciar componentes de interface, permitindo o uso de \texttt{pytest} diretamente
sobre \texttt{models.py}. \\ \hline
RNF03 & \textbf{Confiabilidade}: os resultados numéricos produzidos pelo sistema Python
devem ser equivalentes aos do sistema Java legado, verificados por testes de caixa-preta
comparativos. \\ \hline
RNF04 & \textbf{Portabilidade}: o sistema deve ser executável via navegador
(\texttt{streamlit run app.py}), sem necessidade de JVM, arquivos \texttt{.class} ou
instalação local de dependências além do ambiente Python. \\ \hline
RNF05 & \textbf{Usabilidade}: a interface deve guiar o usuário em um fluxo sequencial e
irreversível (Configuração~$\rightarrow$~Resultados~$\rightarrow$~Gráficos), impedindo
execuções com arquivos incompletos. \\ \hline
RNF06 & \textbf{Desempenho}: o tempo de execução de cada modelo (RLM, \textit{Bagging},
\textit{Subagging}) deve ser medido em milissegundos e reportado como métrica
independente das métricas de acurácia. \\ \hline
RNF07 & \textbf{Escalabilidade}: o sistema deve manter corretude numérica ao processar
matrizes espectrais de grande porte, verificada por testes de estresse sobre matrizes de
grandes dimensões. \\ \hline
RNF08 & \textbf{Internacionalização}: a interface e os componentes visuais devem suportar
múltiplos idiomas (português e inglês) sem duplicação de código, por meio de dicionários
de tradução centralizados. \\ \hline
\end{tabularx}
\begin{tablenotes}[flushleft]
\footnotesize
\item Fonte:~\oautor
\end{tablenotes}
\end{quadro}

Os requisitos funcionais RF03 a RF07 correspondem diretamente às responsabilidades
isoladas em \texttt{models.py}, ao passo que RF09 e RF11 correspondem a
\texttt{components.py}; RF01, RF08, RF10 e RF12 são responsabilidade exclusiva de
\texttt{app.py}. Essa correspondência unívoca entre requisito e módulo confirma, em nível
de especificação, a separação de camadas discutida na Seção~\ref{sec:contratos}.

Os requisitos não funcionais RNF01 a RNF03, por sua vez, são verificados empiricamente
pela suíte de 92 testes automatizados (\texttt{test\_suite.py}) — 21 unitários, 60 de
caixa-preta e 11 de estresse —, detalhada no capítulo de Testes e Validação deste TCC~II.

% ─────────────────────────────────────────────────────────────────────────────
\section{Estratégias de Migração de Sistemas}

A transição do modelo Desktop (Java/Vaadin) para o modelo Web (Python/Streamlit) foi
conduzida por \textbf{reescrita funcional} (\textit{black-box rewrite}), estratégia que
extrai os requisitos do comportamento observável do sistema legado sem reaproveitar o
código-fonte~\cite{becker2020}.

\begin{description}

\item[a)] \textbf{Reescrita Black-Box}: O comentário de cabeçalho do \texttt{app.py}
(\textit{``Conversão do projeto Vaadin Flow (Java) para Streamlit (Python)''}) e as
referências internas ao \texttt{RegressionService.java} indicam que a migração partiu do
comportamento esperado, não do código herdado. Não há dependência de JVM, arquivos
\texttt{.class} ou \texttt{.jar}. O núcleo matemático foi reimplementado diretamente em
\texttt{models.py} com \texttt{numpy}, preservando a equivalência algébrica da
pseudoinversa de Moore-Penrose e dos algoritmos de amostragem.

\item[b)] \textbf{Extração Seletiva da Lógica de Negócio}: A migração isolou apenas o
que tinha valor científico reutilizável: a pseudoinversa (\texttt{perform\_rlm}), a
amostragem com e sem reposição (\texttt{perform\_bagging}), o grid search por validação
cruzada (\texttt{find\_optimal\_parameters}) e as métricas de erro
(\texttt{calc\_metrics}). Camadas de UI do Java — componentes Vaadin, listeners de
evento, gerenciamento de sessão HTTP — foram descartadas e substituídas pelos widgets
declarativos do Streamlit no \texttt{app.py}.

\item[c)] \textbf{Deploy Web Nativo}: O modelo Desktop foi convertido para uma aplicação
\textit{stateless}, executável com \texttt{streamlit run app.py}, containerizável com
Docker e acessível via navegador sem instalação local. Isso reduz o atrito de distribuição
em laboratórios e viabiliza uso colaborativo simultâneo, funcionalidade inexistente no
sistema Vaadin original.

\item[d)] \textbf{Papel e Limites do Estágio Intermediário (Vaadin)}: O estágio Vaadin
Flow funcionou como prova de conceito de exposição web da lógica reaproveitada do Java,
validando a viabilidade de servir o núcleo algorítmico via navegador antes da reescrita
completa. Diferentemente do estágio final em Python, esse estágio não foi acompanhado por
suíte de testes automatizados formal, lacuna reconhecida que motivou a decisão de
formalizar o oráculo de equivalência --- os 92 testes de caixa-preta --- apenas a partir
da reescrita em Streamlit.

\end{description}

% ─────────────────────────────────────────────────────────────────────────────
\section{Mudança de Paradigma Arquitetural}

A reestruturação em três arquivos evidencia uma mudança profunda no paradigma de execução
e no modelo de estado da aplicação.

\begin{description}

\item[a)] \textbf{MVC Tradicional (Java/Vaadin)}: O sistema legado seguia a separação
explícita entre \textit{Model} (dados e regressão), \textit{View} (componentes HTML
gerados pelo Vaadin) e \textit{Controller} (listeners de evento Java). O estado vivia na
sessão HTTP e as atualizações de interface eram parciais e assíncronas (AJAX/push),
permitindo que apenas o componente afetado fosse re-renderizado.

\item[b)] \textbf{Arquitetura em Camadas com Streamlit}: A divisão em
\texttt{models.py} (lógica pura), \texttt{components.py} (fábrica visual) e
\texttt{app.py} (orquestrador) aplica o princípio de \textbf{separação de
responsabilidades} de Sommerville~\cite{sommerville2011} dentro do paradigma Streamlit.
O \texttt{app.py} é um script linear que reexecuta do início ao fim a cada interação do
usuário — clique, upload ou navegação entre abas. Não há controlador permanente; o fluxo
segue \textit{Entrada do Usuário $\rightarrow$ Rerun do Script $\rightarrow$ Renderização
Declarativa}.

\item[c)] \textbf{Gerenciamento de Estado Explícito}: No paradigma Streamlit, o estado
não persiste automaticamente entre reruns. O dicionário
\texttt{st.session\_state["result"]} armazena o output completo de
\texttt{run\_full\_analysis()} — métricas, predições e log — garantindo que as abas
\textbf{Resultados} e \textbf{Gráficos} leiam os dados sem reacionar a regressão. Essa
explicitação do estado é uma diferença fundamental em relação ao modelo MVC com sessão
HTTP implícita do Java.

\end{description}


% =============================================================================
\chapter{Ecossistema de Desenvolvimento (Python e Streamlit)}

\section{Python na Computação Científica}

A escolha do Python como linguagem base é justificada pelo ecossistema maduro voltado à
álgebra linear, ao processamento de sinais espectrais e à quimiometria computacional.
Bibliotecas como \texttt{NumPy}, \texttt{Pandas} e \texttt{SciPy} formam a base
científica sobre a qual o sistema foi construído e pode evoluir, tornando Python a
escolha natural para domínios que lidam com dados espectrométricos de alta
dimensionalidade.

\begin{description}

\item[a)] \textbf{Processamento Vetorial com \texttt{NumPy}}: O \texttt{models.py}
utiliza exclusivamente operações vetorizadas do NumPy:
\texttt{np.linalg.pinv(Xcal)~@~Ycal} para a pseudoinversa, \texttt{np.random.default\_rng}
para amostragem reprodutível com semente configurável, e \texttt{np.sqrt(np.mean(err**2))}
para RMSE. Essas operações substituem laços explícitos em Java, executando em tempo próximo
a linguagens compiladas mesmo para matrizes espectrais com milhares de variáveis.

\item[b)] \textbf{Parsing Robusto com \texttt{parse\_matrix()} e \texttt{parse\_vector()}}:
As funções de leitura em \texttt{models.py} utilizam decodificação UTF-8 com tolerância a
erros, normalização de separador decimal (vírgula para ponto) e filtragem de linhas
vazias. Esse parsing manual garante compatibilidade com arquivos \texttt{.txt} gerados
por instrumentos espectrométricos sem dependência de formatação rígida.

\item[c)] \textbf{Interoperabilidade Científica com \texttt{Pandas}}: A integração com
\texttt{pandas} no \texttt{app.py} permite exportar a tabela comparativa de métricas
diretamente em CSV (\texttt{df\_metrics.to\_csv()}), mantendo o fluxo de trabalho típico
de laboratórios de quimiometria. O \texttt{pd.DataFrame} estrutura os resultados de RMSE,
MAE, R² e tempo de execução de forma tabular, viabilizando análise posterior em planilhas
ou pipelines de dados sem conversões intermediárias.

\item[d)] \textbf{Extensão Natural para \texttt{SciPy} em Sinais Espectrométricos}: O
\texttt{SciPy} é a extensão natural do ecossistema para o domínio de espectrometria.
Funções como \texttt{scipy.signal.savgol\_filter()} realizam suavização de Savitzky--Golay
sobre espectros ruidosos; \texttt{scipy.fft} calcula transformadas para análise no domínio
da frequência; e \texttt{scipy.signal.find\_peaks()} identifica automaticamente picos em
espectros de massa (eixo $m/z$ vs.\ Intensidade). A arquitetura modular do
\texttt{models.py} — funções puras recebendo e retornando \texttt{np.ndarray} — permite
incorporar essas rotinas de pré-processamento espectral sem alterações na camada de
apresentação, tornando a adoção de \texttt{SciPy} uma evolução natural e de baixo risco
para versões futuras do sistema.

\end{description}

% ─────────────────────────────────────────────────────────────────────────────
\section{Ciclo de Vida do Streamlit}

O Streamlit opera sob um modelo de execução declarativo e \textit{stateless} por padrão,
com implicações diretas no tempo de resposta e na arquitetura do \texttt{app.py}.

\begin{description}

\item[a)] \textbf{Execução Top-Down Unidirecional}: A cada interação do usuário — clique
em botão, alteração de campo de texto ou navegação entre abas — o interpretador Python
relê o \texttt{app.py} do início ao fim. Não há um controller permanente em memória.
O fluxo segue estritamente \textit{Entrada $\rightarrow$ Rerun $\rightarrow$ Renderização},
tornando o comportamento da aplicação previsível e auditável, mas exigindo atenção ao
custo computacional de cada instrução no nível superior do script.

\item[b)] \textbf{Persistência via \texttt{st.session\_state}}: Para evitar que o
\texttt{run\_full\_analysis()} seja re-executado a cada rerun, o \texttt{app.py}
armazena o resultado completo em \texttt{st.session\_state["result"]}. As abas
\textbf{Resultados} e \textbf{Gráficos} verificam a existência dessa chave antes de
renderizar, lendo os dados diretamente da memória de sessão sem acionar novamente os
laços de otimização.

\item[c)] \textbf{Feedback Assíncrono Simulado}: A barra de progresso
(\texttt{st.progress}) é atualizada via callback \texttt{update\_progress} passado para
\texttt{run\_full\_analysis()}, que o propaga internamente a
\texttt{find\_optimal\_parameters()}. Embora o Streamlit seja síncrono por natureza, essa
estratégia de callback mitiga a percepção de travamento durante as 150 execuções do grid
search (5 valores de \textit{m} + 5 $\times$ 5 combinações de \textit{k} $\times$
\textit{m}, cada uma com 5 folds).

\end{description}

% ─────────────────────────────────────────────────────────────────────────────
\section{Gerenciamento de Memória e Performance}

A gestão eficiente de recursos é crítica para manter a responsividade durante o
processamento de matrizes espectrais de alta dimensionalidade.

\begin{description}

\item[a)] \textbf{Uso Estratégico de \texttt{st.session\_state}}: Com a separação em
\texttt{models.py} e \texttt{app.py}, o estado mantém o dicionário de saída de
\texttt{run\_full\_analysis()} vivo entre reruns. Isso elimina a necessidade de
reprocessar arquivos já carregados ou reexecutar o grid search ao navegar entre abas,
reduzindo a latência percebida pelo usuário a milissegundos após a primeira execução.

\item[b)] \textbf{Oportunidade de \texttt{@st.cache\_data}}: Como \texttt{models.py}
contém funções puras — \texttt{parse\_matrix()}, \texttt{parse\_vector()} e
\texttt{run\_full\_analysis()} recebem apenas dados serializáveis e retornam resultados
determinísticos — o decorador \texttt{@st.cache\_data} pode ser aplicado diretamente,
sem modificações na lógica. Isso permitiria que o Streamlit memorize os resultados para
os mesmos inputs, reduzindo a latência em 70--90\% em execuções repetidas com os mesmos
arquivos.

\item[c)] \textbf{Isolamento da Carga Computacional}: As funções pesadas
(\texttt{find\_optimal\_parameters}, \texttt{perform\_bagging}) estão encapsuladas em
\texttt{models.py}, recebem \texttt{np.ndarray} e retornam resultados determinísticos.
Esse isolamento evita que operações de UI do \texttt{app.py} interfiram nos cálculos e
facilita a futura migração para processamento assíncrono em background via
\texttt{concurrent.futures} ou \texttt{asyncio}, sem alterações na camada de
apresentação.

\end{description}

% ─────────────────────────────────────────────────────────────────────────────
\section{IDEs e Ambiente de Desenvolvimento}

O desenvolvimento e a manutenção dos três módulos são potencializados pelo VS~Code
(\textit{Visual Studio Code}), editor de código aberto da Microsoft que se consolidou
como padrão da indústria para projetos Python científicos~\cite{vscode2024}.

\begin{description}

\item[a)] \textbf{Extensões Especializadas para o Projeto}: A extensão \texttt{Python}
(Microsoft) provê IntelliSense e navegação entre os três módulos; \texttt{Pylance} realiza
\textit{type checking} estático, detectando incompatibilidades de \texttt{np.ndarray}
entre \texttt{models.py} e \texttt{app.py} antes da execução. A extensão
\texttt{Jupyter} permite prototipar e validar funções matemáticas isoladamente em
\textit{notebooks} antes de integrá-las ao \texttt{models.py}, acelerando o ciclo de
desenvolvimento.

\item[b)] \textbf{Qualidade e Padronização de Código}: A integração nativa com
\texttt{black} (formatação automática), \texttt{ruff} ou \texttt{flake8} (linting de
estilo e erros) e \texttt{mypy} (validação de tipos) garante consistência entre os três
arquivos. Os \textit{type hints} adotados em \texttt{models.py} e \texttt{components.py}
são validados em tempo real, prevenindo chamadas com tipos incorretos a partir do
\texttt{app.py}.

\item[c)] \textbf{Depuração com Breakpoints Cirúrgicos}: O arquivo \texttt{launch.json}
do VS~Code permite inserir breakpoints diretamente dentro de
\texttt{run\_cross\_validation()} ou \texttt{calc\_metrics()}, inspecionando o
\textit{shape} dos arrays, os valores intermediários de \texttt{rmse\_list} e o índice
do fold atual — informações impossíveis de obter com \texttt{st.write()} na interface.
O terminal integrado gerencia ambientes isolados via \texttt{venv} ou \texttt{conda},
garantindo que as dependências (\texttt{streamlit}, \texttt{numpy}, \texttt{plotly},
\texttt{pandas}) não conflitem com outros projetos do sistema operacional.

\end{description}


% =============================================================================
\chapter{Interface e Experiência do Usuário (UI/UX) Científica}

\section{Design Centrado em Dados}

A interface do \texttt{app.py} organiza o fluxo analítico de forma hierárquica,
priorizando a clareza e a produtividade do cientista em detrimento de elementos
decorativos. A distribuição dos elementos entre abas, colunas e área central segue
princípios de design centrado no usuário, com decisões arquiteturais explícitas sobre
onde cada tipo de conteúdo deve residir.

\begin{description}

\item[a)] \textbf{Navegação por Etapas (\texttt{st.tabs})}: As três abas —
\textbf{Configuração}, \textbf{Resultados} e \textbf{Gráficos} — guiam o usuário em um
fluxo linear e irreversível: entrada de dados $\rightarrow$ execução $\rightarrow$
inspeção de resultados. Essa estrutura espelha o \textit{workflow} típico de validação de
modelos quimiométricos e reduz a sobrecarga cognitiva ao esconder informações irrelevantes
para a etapa atual.

\item[b)] \textbf{Layout em Colunas e Hierarquia Visual}: O \texttt{app.py} utiliza
\texttt{st.columns(2)} para agrupar uploads de calibração e teste lado a lado,
sinalizando visualmente a paridade entre os conjuntos. A área central é reservada para
cards de métricas gerados por \texttt{components.metric\_card\_html()} e para a tabela
comparativa em \texttt{pd.DataFrame}. O CSS centralizado em \texttt{components.APP\_CSS}
— incluindo gradientes de fundo, bordas coloridas por modelo e estilo terminal para o
log — replica dashboards técnicos e industriais sem poluir o \texttt{app.py} com
definições de estilo.

\item[c)] \textbf{Feedback Progressivo}: O botão \textbf{Executar Análise} é
desabilitado (\texttt{disabled=not files\_ok}) enquanto algum dos quatro arquivos estiver
ausente, prevenindo execuções inválidas. A barra de progresso atualizada pelo callback e
a mensagem \texttt{st.success} ao término fornecem confirmação clara de conclusão,
padrão recomendado pelas diretrizes de usabilidade para interfaces científicas.

\item[d)] \textbf{Decisão Arquitetural: \texttt{st.tabs} vs.\ \texttt{st.sidebar}}: O
Streamlit oferece dois paradigmas de organização de parâmetros técnicos e resultados:
a \texttt{st.sidebar} (painel lateral persistente, visível em todas as telas) e o
\texttt{st.tabs} (abas sequenciais que segmentam o fluxo). Para o perfil de usuário
científico deste sistema — que opera em etapas ordenadas e não retorna aos parâmetros
após a execução — o modelo de abas foi preferido por três razões: (1) impõe
explicitamente a sequência \textit{Configuração $\rightarrow$ Resultados $\rightarrow$
Gráficos}, prevenindo execuções prematuras; (2) preserva a área central integralmente
para visualização de resultados e gráficos, sem que uma barra lateral compita com os
gráficos do Plotly em largura de tela; e (3) alinha-se ao padrão de \textit{wizards}
adotado em \textit{softwares} de análise científica (Origin, MATLAB App Designer), nos
quais cada etapa revela seu conteúdo apenas quando os pré-requisitos anteriores foram
cumpridos. A \texttt{st.sidebar} permanece como alternativa recomendada para versões
futuras que implementem parâmetros globais persistentes, como configuração de idioma,
normalização de dados ou seleção de paleta de cores para publicação.

\end{description}

% ─────────────────────────────────────────────────────────────────────────────
\section{Visualização Interativa com Plotly}

A escolha do Plotly no \texttt{components.py} é fundamentada pela necessidade de
inspeção técnica precisa em espectrometria e regressão multimodelo.

\begin{description}

\item[a)] \textbf{Integração Nativa e Centralizada}: Com a separação em
\texttt{components.py}, todas as configurações visuais — cores por modelo
(\texttt{COR\_RLM}, \texttt{COR\_BAGGING}, \texttt{COR\_SUBAGGING}), estilos de eixo
(\texttt{\_AXIS\_STYLE}), cores de fundo (\texttt{\_PLOT\_BG}) — estão definidas em um
único arquivo. Alterar a paleta de cores ou o estilo de grade para padrões de publicação
científica requer modificação em um só lugar, sem tocar no \texttt{app.py} ou no
\texttt{models.py}.

\item[b)] \textbf{Ferramentas de Inspeção Científica}: Os gráficos \texttt{go.Scatter}
gerados por \texttt{build\_chart()} e \texttt{build\_scatter\_chart()} habilitam zoom,
\textit{pan}, \textit{hover} com coordenadas exatas e legendas interativas. A linha ideal
\textit{y=x} (\texttt{mode="lines", dash="dot"}) no gráfico de dispersão
(\texttt{build\_scatter\_chart}) permite validação visual imediata de viés e dispersão
dos três modelos sobrepostos — RLM Simples, Bagging e Subagging — em uma única figura.

\item[c)] \textbf{Padrões de Publicação}: Os eixos customizados (\texttt{showgrid=True},
\texttt{gridcolor="\#000000"}, fontes pretas explícitas via \texttt{\_AXIS\_STYLE})
garantem legibilidade para relatórios técnicos, apresentações e submissão a periódicos,
alinhando-se às normas de visualização de dados científicos sem exigir pós-processamento
externo das figuras.

\end{description}

% ─────────────────────────────────────────────────────────────────────────────
\section{Contratos de Interface e Funções}
\label{sec:contratos}

A separação em três módulos aplica rigorosamente o princípio de separação de
responsabilidades, tornando o sistema testável, manutenível e extensível
independentemente em cada camada.

\begin{description}

\item[a)] \textbf{Separação de Camadas}: O \texttt{models.py} não importa
\texttt{streamlit} nem \texttt{plotly}. O \texttt{components.py} não importa
\texttt{streamlit} nem \texttt{numpy}. O \texttt{app.py} não contém nenhuma função
matemática — a instrução mais complexa é a chamada \texttt{result =
run\_full\_analysis(...)}. Essa separação permite testar o \texttt{models.py}
isoladamente com \texttt{pytest} e dados sintéticos, sem instanciar nenhum componente de
interface.

\item[b)] \textbf{Contratos Explícitos via \textit{Type Hints}}: Todas as funções do
\texttt{models.py} declaram tipos de entrada e saída:
\texttt{perform\_bagging(...) -> np.ndarray},
\texttt{calc\_metrics(...) -> tuple[float, float, float]},
\texttt{find\_optimal\_parameters(...) -> tuple[dict, str]}.
No \texttt{app.py}, \texttt{parse\_int\_list()} e \texttt{parse\_float\_list()} lançam
\texttt{ValueError} com mensagens orientativas para entradas fora do domínio válido —
inteiros positivos para \textit{m} e decimais em $(0,1]$ para \textit{k} — prevenindo
falhas silenciosas na camada de negócio.

\item[c)] \textbf{Manutenibilidade e Escalabilidade}: Adicionar um novo modelo preditivo
(como Ridge Regression ou Random Forest) requer apenas uma nova função em
\texttt{models.py} e um novo bloco de exibição em \texttt{app.py}. Modificar a paleta de
cores ou o layout dos gráficos requer apenas alterações em \texttt{components.py}.
Nenhuma dessas mudanças impacta as outras camadas, confirmando que a arquitetura modular
adotada atende ao critério de \textbf{abertura para extensão, fechamento para
modificação} (\textit{Open/Closed Principle}) proposto por Robert C.
Martin~\cite{martin2008}.

\end{description}
